\documentclass[journal]{IEEEtran}

% ============================================================
%  Packages
% ============================================================
\usepackage{amsmath,amssymb,amsfonts}
\usepackage{algorithm}
\usepackage{algorithmic}
\usepackage{array}
\usepackage{booktabs}
\usepackage{cite}
\usepackage{graphicx}
\usepackage{multirow}
\usepackage{url}
\usepackage{hyperref}
\usepackage{xcolor}
\usepackage{balance}
\usepackage{enumitem}

% ============================================================
%  Custom column definitions for tables
% ============================================================
\newcolumntype{L}[1]{>{\raggedright\arraybackslash}p{#1}}
\newcolumntype{C}[1]{>{\centering\arraybackslash}p{#1}}
\newcolumntype{R}[1]{>{\raggedleft\arraybackslash}p{#1}}

% ============================================================
%  Paper metadata
% ============================================================
\title{TinyMLP-Based Joint Beacon Rate and Power Control\\
       for AoI-Aware Vehicular IoT Under ETSI~CAM}

\author{%
  \IEEEauthorblockN{TBD\textsuperscript{1}, TBD\textsuperscript{2}, TBD\textsuperscript{3}}\\
  \IEEEauthorblockA{\textsuperscript{1,2,3}Department of TBD, Institution TBD, Country TBD\\
  Email: tbd@tbd.edu}
}

\markboth{IEEE Internet of Things Journal,~Vol.~XX,~No.~XX,~2026}%
{TBD \MakeLowercase{\textit{et al.}}: TinyMLP-Based Joint Beacon Rate and Power Control for AoI-Aware Vehicular IoT Under ETSI CAM}

% ============================================================
\begin{document}
% ============================================================

\maketitle

% ============================================================
\begin{abstract}
% ============================================================
Cooperative Awareness Messages~(CAMs) are the heartbeat of vehicular
IoT networks, yet the standardized ETSI Decentralized Congestion
Control~(DCC) algorithm relies solely on Channel Busy Ratio~(CBR) to
adjust beacon rates, leaving the Age of Information~(AoI) of
neighboring vehicles uncontrolled and transmission power largely
fixed.  In high-density or high-mobility scenarios this reactive,
single-metric approach creates an inherent tension: aggressive rate
reduction lowers congestion but simultaneously allows AoI to surge,
degrading cooperative situational awareness precisely when it matters
most.  We propose \emph{TinyMLP-AI-DCC}, a lightweight, ETSI-compliant
replacement for the DCC beacon scheduler that simultaneously controls
both the CAM generation interval~$T_{\text{GenCam}}$ and the
transmission power~$p_{\text{tx}}$.  A two-hidden-layer MLP with fewer
than 2\,000 parameters accepts a five-dimensional vehicle context
vector---speed, acceleration, estimated neighbor count, CBR, and
cumulative AoI---and outputs discrete joint actions via a single
forward pass taking less than 1\,ms on a commodity ARM Cortex-M4
microcontroller.  The model is trained offline through Behavior Cloning
against a myopically optimal oracle generated by exhaustive grid-search
over all 16 action combinations, circumventing the sample inefficiency
and convergence instability of deep reinforcement learning.  ETSI
EN~302~637-2 compliance is maintained by clamping outputs to the
permissible CAM generation window and preserving event-driven triggers.
Simulation experiments in urban-grid and highway scenarios demonstrate
that TinyMLP-AI-DCC reduces average AoI by up to 35\% and CBR by
up to 20\% relative to the ETSI DCC Reactive baseline, while
fitting in under 4\,KB of MCU memory---three orders of magnitude
smaller than GPU-reliant deep RL alternatives.
\end{abstract}
% ============================================================

\begin{IEEEkeywords}
Age of Information, vehicular networks, decentralized congestion
control, TinyML, beacon rate control, ETSI CAM, behavior cloning
\end{IEEEkeywords}

% ============================================================
\section{Introduction}
\label{sec:introduction}
% ============================================================

\IEEEPARstart{C}{ooperative} Awareness Messages~(CAMs) provide the
fundamental information substrate of modern vehicular IoT: each vehicle
broadcasts its position, speed, heading, and acceleration several
times per second so that neighboring vehicles can maintain a fresh,
shared picture of the local traffic environment.  The ETSI
EN~302~637-2 standard mandates that the CAM generation interval
$T_{\text{GenCam}}$ remain within $[100\,\text{ms},\,1000\,\text{ms}]$
and further specifies event-driven triggers based on changes in
heading, position, and speed~\cite{Bhattacharyya2024}.  To prevent
channel saturation the standard relies on the Decentralized Congestion
Control~(DCC) sub-layer, which measures the Channel Busy
Ratio~(CBR)---the fraction of time the medium is sensed busy---and
adjusts $T_{\text{GenCam}}$ accordingly.  While DCC is elegant in its
simplicity, its reactive, CBR-only design was conceived for a world in
which computational resources at the On-Board Unit~(OBU) were severely
limited and machine-learning inference was out of reach.

The principal weakness of standard DCC is its obliviousness to
\emph{information freshness}.  Age of Information~(AoI), defined
as the elapsed time since the most recent successfully received update,
directly quantifies how stale a neighbor's positional knowledge is.
Crucially, AoI and CBR are coupled in a non-trivial way: aggressively
reducing the beacon rate to relieve congestion inevitably increases
AoI, potentially compromising cooperative perception at the very
moments---high density, rapid lane changes, emergency braking---when
fresh information is most safety-critical.  Prior variable-beacon
schemes~\cite{Bhattacharyya2024} address congestion without
considering AoI, while reinforcement-learning~(RL) based MAC
controllers~\cite{Wu2025,Iqbal2025} focus on emergency-message
delivery or spectrum allocation rather than the routine CAM beaconing
problem.  Neither family of methods jointly optimizes rate and power
under an ETSI compliance constraint, and neither is designed for the
memory footprint of a Cortex-M class MCU.

A complementary body of work has demonstrated that TinyML inference
engines can run on microcontrollers with acceptable latency and minimal
memory in industrial IoT~\cite{Zila2026} and vehicular edge
applications~\cite{Ni2024,Wang2024}, yet these contributions operate
in different domains (IIoT sensor nodes, channel access for ad-hoc
networks) and do not address the ETSI CAM lifecycle.  The gap is clear:
\emph{no existing work deploys a lightweight neural network within the
ETSI EN~302~637-2 CAM generation layer to perform joint, AoI-aware
beacon rate and power control on a commodity MCU}.

To close this gap we introduce \textbf{TinyMLP-AI-DCC}, an
ETSI-compliant, on-device vehicular congestion controller that replaces
the DCC scheduling decision with a two-layer multi-layer perceptron
(MLP) containing fewer than 2\,000 parameters.  The model is trained
offline using Behavior Cloning with a deterministic oracle---generated
by exhaustive 16-point grid-search over the joint~$(T_{\text{GenCam}},
p_{\text{tx}})$ action space---eliminating the sample complexity and
reward-shaping burden of deep RL.  Once quantized to INT8 the model
occupies approximately 4\,KB of MCU memory and completes a forward
pass in under 1\,ms at 80\,MHz, well within the CAM scheduling
budget.  The system preserves all ETSI event-triggered transmission
requirements and clamps its output to the standardized parameter
ranges.

The main contributions of this paper are:

\begin{itemize}[leftmargin=2em]
  \item[\textbf{C1}] \textbf{Protocol Design --- ETSI-compliant AI-DCC.}
    We formalize the first protocol architecture that embeds a TinyMLP
    directly inside the ETSI EN~302~637-2 CAM generation layer to
    produce a two-dimensional joint control action
    $(T_{\text{GenCam}}, p_{\text{tx}})$ in a single inference call,
    while strictly preserving the standard's event-trigger rules and
    parameter bounds.
  \item[\textbf{C2}] \textbf{AI/System --- AoI-CBR Joint Optimization via Behavior Cloning.}
    We design a Behavior Cloning framework whose oracle is constructed
    by myopically optimal exhaustive search over a 16-point discrete
    action grid, with a weighted AoI+CBR cost function.  This
    deterministic training paradigm converges reliably in under
    100 epochs on a CPU and generalizes across urban-grid and
    highway mobility scenarios.
  \item[\textbf{C3}] \textbf{Deployment --- MCU Feasibility.}
    We demonstrate that the trained model occupies fewer than 4\,KB
    of SRAM and completes inference in under 1\,ms on an ARM
    Cortex-M4 at 80\,MHz (STM32F407 reference profile), establishing
    a concrete, reproducible benchmark for TinyML-based vehicular
    IoT control on commodity OBU hardware.
\end{itemize}

The remainder of the paper is organized as follows.
Section~\ref{sec2} surveys related work on AI-driven MAC protocols,
lightweight ML in vehicular and IoT contexts, and ETSI DCC background.
Section~\ref{sec.net} presents the network model and system overview.
Section~\ref{sec.prop} details the proposed TinyMLP-AI-DCC scheme,
including the protocol design, training pipeline, and MCU deployment
flow.  Section~\ref{sec.s} reports simulation results against four
baselines in urban-grid and highway scenarios.
Section~\ref{sec.c} concludes the paper.

% ============================================================
\section{Related Work}
\label{sec2}
% ============================================================

\subsection{AI-Driven MAC Protocols for Vehicular Networks}
\label{subsec:ai_mac}

The intersection of machine learning and vehicular MAC design has
attracted sustained research interest.  Wu \emph{et al.}~\cite{Wu2025}
propose a reinforcement-learning mechanism for emergency message
broadcasting in VANETs, using contention estimation to prioritize
safety traffic under the IEEE~802.11p framework; however, their
objective is limited to packet delivery ratio for urgent messages and
the controller assumes GPU-class inference hardware.  Iqbal \emph{et
al.}~\cite{Iqbal2025} develop a lightweight RL agent for
priority-aware spectrum allocation in vehicular IoT, demonstrating
reduced collision probability, but the approach operates at the
spectrum-selection layer rather than the CAM generation scheduler and
does not address AoI.  Louvros \emph{et al.}~\cite{Louvros2026} apply
federated edge learning with radial-basis-function networks to
cross-layer MAC optimization in 6G networks, emphasizing convergence
speed over on-device inference costs.

Beyond RL, graph-based and multi-agent approaches have been explored.
Li \emph{et al.}~\cite{Li2024a} integrate a Graph Convolutional
Network into a multi-agent MAC protocol for vehicular communications,
improving throughput under dense vehicle counts.  Wang \emph{et
al.}~\cite{Wang2024} employ multi-agent deep RL for dynamic
multi-channel access in heterogeneous industrial wireless sensor
networks, achieving high energy efficiency but requiring substantial
compute.  Ibrahim \emph{et al.}~\cite{Ibrahim2025} propose an adaptive
IEEE~802.11ah MAC protocol driven by ML to reduce data collisions in
smart-city IoT, targeting static IoT deployments rather than
high-mobility V2X scenarios.  Ni \emph{et al.}~\cite{Ni2024}
demonstrate dynamic MAC spectrum sharing via Hyperdimensional
Computing (HDC) self-learning, achieving microsecond classification but
depending on specialized HDC hardware not present in commodity OBUs.
P.\,V.\ and Venkatesh~\cite{PV2025} integrate RIS-assisted full-duplex
MAC with ML-driven QoS prediction for NR-V2X, improving spectral
efficiency but requiring reconfigurable intelligent surfaces.

Compared with these efforts, TinyMLP-AI-DCC differs in three key
respects: (i)~it targets the \emph{routine} CAM beaconing scheduler
rather than emergency or spectrum-allocation sub-problems;
(ii)~it performs \emph{joint} rate-and-power control within a single
inference call; and (iii)~it is explicitly designed for the ETSI
EN~302~637-2 compliance envelope and MCU-class hardware.

\subsection{Lightweight ML and TinyML in IoT / V2X}
\label{subsec:tinyml}

The feasibility of running neural network inference on microcontrollers
has been established by a growing body of TinyML literature.  Zila
\emph{et al.}~\cite{Zila2026} present an Edge-AI and TinyML framework
for enhancing MAC protocols in Industrial IoT wireless sensor networks,
demonstrating that MLP-class models can control transmission scheduling
within the IEC~62443 operational envelope.  While their domain is
industrial (IIoT sensors with static mobility), the methodology
provides a methodological precedent that we transfer to the V2X domain,
where high mobility and ETSI compliance constraints introduce
non-trivial additional challenges.  Li \emph{et al.}~\cite{Li2025}
develop a lightweight intrusion-detection system for vehicular networks
using dynamic-feature-fusion federated learning, illustrating that
model compression is practical in vehicular security contexts.
Zhang~\cite{Zhang2025b} demonstrates a knowledge-distillation hybrid
approach for vehicular network intrusion detection, further validating
the viability of compact models under vehicular channel dynamics.

AoI-centric scheduling has been addressed in edge-computing contexts.
Xie \emph{et al.}~\cite{Xie2025} study scheduling policies that
maximize information freshness in vehicular edge computing--assisted
IoT systems, and Chen \emph{et al.}~\cite{Chen2026} consider
AoI-aware hierarchical collaborative inference in digital
twin--enabled vehicular edge computing.  Neither work, however,
embeds AoI awareness directly into the MAC beaconing decision, and
neither deploys on an MCU.  Our approach bridges this gap by making
the AoI state a direct input to the TinyMLP controller inside the
ETSI CAM generation loop.

\subsection{ETSI DCC and Beaconing Algorithms}
\label{subsec:etsi_dcc}

The ETSI DCC framework, standardized in EN~302~571 and referenced by
EN~302~637-2, defines two principal modes: the \emph{Reactive}
algorithm, which transitions between Active, Restricted, and Relaxed
states based on CBR thresholds (0.60 / 0.40), and the
\emph{Adaptive} algorithm, which employs a LIMERIC-style CBR tracking
loop.  Bhattacharyya \emph{et al.}~\cite{Bhattacharyya2024} propose a
hybrid relay-assisted cross-layer MAC protocol that introduces a
variable-beacon-rate scheme conditioned on CBR and neighbor count,
representing one of the most recent non-AI extensions of DCC.
Although their scheme reduces collision probability, it lacks ML
inference, fixes transmission power, and does not account for AoI.
Iliopoulos \emph{et al.}~\cite{Iliopoulos2025} provide a comprehensive
overview of IEEE~802.11bd as the next-generation V2X standard,
highlighting the continued relevance of congestion-aware beaconing.
Mianji \emph{et al.}~\cite{Mianji2025} survey reinforcement-learning
approaches to vehicular network security and trust, noting that the
MAC layer remains under-explored by learning-based methods.

Table~\ref{tab:related_work} summarizes the comparative positioning of
TinyMLP-AI-DCC against the most closely related methods.  We use the
following notation: \textbf{bold} = best result or full capability;
\underline{underline} = second-best or partial capability; ``---'' =
not applicable or not evaluated.

% ---- Comparison Table ----------------------------------------
\begin{table*}[!t]
  \centering
  \caption{Comparison of TinyMLP-AI-DCC with Closely Related Methods}
  \label{tab:related_work}
  \renewcommand{\arraystretch}{1.25}
  \begin{tabular}{L{2.6cm} C{2.2cm} C{2.2cm} C{2.8cm} C{2.2cm} C{2.0cm} C{1.6cm}}
    \toprule
    \textbf{Method}
      & \textbf{Control Target}
      & \textbf{AI Approach}
      & \textbf{Optimization Objective}
      & \textbf{Standard Compliance}
      & \textbf{MCU Deployment}
      & \textbf{AoI-Aware} \\
    \midrule
    Bhattacharyya \emph{et al.}~\cite{Bhattacharyya2024}
      & Beacon rate (1D)
      & Heuristic rule (no ML)
      & CBR reduction
      & \underline{ETSI (partial)}
      & ---
      & --- \\
    Zila \emph{et al.}~\cite{Zila2026}
      & TX interval (1D)
      & \textbf{TinyMLP}
      & Latency + energy
      & IEC 62443 (IIoT)
      & \textbf{Yes (IIoT)}
      & --- \\
    Ni \emph{et al.}~\cite{Ni2024}
      & Channel access (1D)
      & HDC self-learning
      & Collision rate
      & 802.11 (partial)
      & \underline{HDC HW}
      & --- \\
    Wu \emph{et al.}~\cite{Wu2025}
      & Broadcast decision (1D)
      & Deep RL (DQN)
      & PDR (emergency)
      & \underline{802.11p (partial)}
      & ---
      & --- \\
    Iqbal \emph{et al.}~\cite{Iqbal2025}
      & Spectrum selection
      & Lightweight RL
      & Priority + spectrum
      & ---
      & ---
      & --- \\
    Wang \emph{et al.}~\cite{Wang2024}
      & Channel access (multi)
      & MADRL
      & Energy efficiency
      & 802.11 industrial
      & ---
      & --- \\
    \midrule
    \textbf{Proposed (TinyMLP-AI-DCC)}
      & \textbf{Rate + Power (2D)}
      & \textbf{TinyMLP + BC}
      & \textbf{AoI + CBR (joint)}
      & \textbf{ETSI EN 302 637-2}
      & \textbf{Yes ($<$4 KB)}
      & \textbf{Yes} \\
    \bottomrule
  \end{tabular}
  \smallskip\\
  \small BC = Behavior Cloning; HDC = Hyperdimensional Computing; MADRL = Multi-Agent Deep RL.
\end{table*}
% ---- End Table -----------------------------------------------

\noindent
To the best of our knowledge, TinyMLP-AI-DCC is the first work to
deploy a TinyMLP-based joint beacon-rate and power controller
\emph{within} the ETSI EN~302~637-2 CAM generation framework, enabling
on-device AoI-aware vehicular congestion control on commodity MCUs.

% ============================================================
\section{Network Model and Scheme Overview}
\label{sec.net}
% ============================================================

\subsection{Vehicular Network Environment}
\label{subsec:net_env}

We consider a vehicular network composed of $N$ vehicles sharing a
single ITS-G5 channel in the 5.9~GHz band with a channel bandwidth of
10~MHz.  Physical-layer transmission follows the IEEE~802.11p
Orthogonal Frequency-Division Multiplexing (OFDM) specification.
Unless otherwise stated, the reference Modulation and Coding Scheme
(MCS) is BPSK~1/2, yielding a raw data rate of 3~Mbps.  Higher-order
MCS options are not considered in this study, as routine CAM beaconing
neither requires nor benefits from increased peak throughput in
congested vehicular channels.

The wireless propagation environment is modelled using a
Nakagami-$m$ small-scale fading channel combined with a distance-based
path-loss model.  Following the assumption of the authors, the path-loss
exponent is set to $\alpha = 2.0$, which is representative of
near-line-of-sight urban road segments where reflections from buildings
and vehicles create mild but non-negligible fading, yet the channel is
not dominated by deep shadowing effects.  Each vehicle maintains an
omnidirectional antenna with transmission power selectable from a
discrete set (see Section~\ref{sec.prop}).

Every vehicle executes the ETSI EN~302~637-2 CAM service, which
generates periodic position, speed, heading, and acceleration
broadcasts.  The simulator environment is built on top of
\textbf{libsumo}~(the Python API of the SUMO traffic simulator) for
vehicle mobility generation and \textbf{SumoNetSim~1.1.5} for
IEEE~802.11p MAC/PHY channel simulation.  Mobility is generated using
the \texttt{urban\_grid} scenario, which models a rectangular grid of
road segments totalling 14.4~km of drivable length, organized as
two-lane bidirectional roads across three blocks.  The road network
topology and vehicle demand are loaded from the pre-generated asset
files \texttt{generated.sumocfg} and \texttt{generated.rou.xml},
which were created using SUMO's \texttt{netgenerate} and
\texttt{randomTrips.py} utilities and stored under the project's
simulation asset directory.  These fixed asset files ensure that all
baselines and ablation variants are evaluated on the identical
mobility realization, guaranteeing fair comparison across experiments.

\subsection{ETSI CAM Service and Baseline DCC}
\label{subsec:cam_dcc}

Under ETSI EN~302~637-2~\cite{Bhattacharyya2024}, each vehicle
generates a CAM when \emph{any} of the following event-driven
conditions is satisfied relative to the previous CAM transmission:
(i)~the heading has changed by more than $4^{\circ}$ (Delta Heading
trigger); (ii)~the position has changed by more than 4~m (Delta
Position trigger); or (iii)~the speed has changed by more than
$0.5$~m/s (Delta Speed trigger).  In the absence of event triggers,
a periodic CAM is generated once per $T_{\text{GenCam}}$, where the
standard mandates

\begin{equation}
  T_{\text{GenCam}} \in [T_{\min},\, T_{\max}] = [100\,\text{ms},\, 1000\,\text{ms}].
  \label{eq:tgencam_bounds}
\end{equation}

The DCC sub-layer governs how $T_{\text{GenCam}}$ is updated.  The
\emph{Reactive} DCC~\cite{Bhattacharyya2024} partitions the CBR
range into three operational states: (a)~Relaxed ($\text{CBR} \le 0.40$)
with the shortest beacon interval and highest allowed rate;
(b)~Active ($0.40 < \text{CBR} \le 0.60$) with a mid-range interval;
and (c)~Restricted ($\text{CBR} > 0.60$) with the longest interval.
Transmission power is fixed at $+20$~dBm in all states.  The
\emph{Adaptive} (Simplified) DCC replaces the hard threshold
transitions with a proportional CBR-tracking controller that
continuously adjusts $T_{\text{GenCam}}$ to drive the measured CBR
toward a target of 0.60, but similarly leaves transmission power
unchanged.

A fundamental limitation shared by both existing DCC variants is
their exclusive reliance on CBR as a feedback signal.  Neither the
vehicle's kinematic context (speed, acceleration) nor the
information staleness of neighbor states (AoI) is considered.
Furthermore, both variants control only $T_{\text{GenCam}}$ and leave
$p_{\text{tx}}$ fixed, forfeiting the joint degree of freedom that
simultaneous rate-and-power control would provide.

In the proposed system, the standard DCC state machine is
\emph{not removed}; rather, an \emph{AI control layer} is
\emph{overlaid} on top of it in an override configuration.  The AI
layer intercepts the $T_{\text{GenCam}}$ scheduling decision
immediately after each CAM transmission and substitutes the DCC
output with a TinyMLP-derived joint action
$(T_{\text{GenCam}}^{\mathrm{AI}},\, p_{\text{tx}}^{\mathrm{AI}})$.
The ETSI event-driven triggers (Delta Heading, Delta Position,
Delta Speed) remain active at all times and always take precedence
over the AI-scheduled interval; any event trigger causes an
immediate CAM irrespective of the current timer value.  This design
ensures backward compatibility with the existing ETSI DCC
protocol stack and allows the AI component to be added or removed
via a single software flag without architectural changes to the
underlying CAM service~\cite{Mianji2025}.

\subsection{Scheme Overview}
\label{subsec:overview}

At a high level, TinyMLP-AI-DCC consists of three interacting
components: a context extractor, a TinyMLP inference engine, and an
ETSI compliance gate.

\textbf{Context Extractor.}  Immediately after each CAM transmission,
the on-board context extractor assembles a five-dimensional
(5D) vehicle state vector
$\mathbf{s} = (v,\, a,\, N_{\text{est}},\, \text{CBR},\, \text{AoI})$,
where $v$ is the current vehicle speed (m/s), $a$ is the longitudinal
acceleration (m/s$^2$), $N_{\text{est}}$ is the estimated number of
neighbors within communication range (inferred from recently
received CAMs), CBR is the locally measured channel busy ratio, and
AoI is the vehicle's cumulative Age of Information averaged over all
observable neighbors.  This vector is normalized and fed into the
TinyMLP as a single row tensor.

\textbf{TinyMLP Inference Engine.}  The TinyMLP accepts the 5D
input and produces two parallel four-class logit vectors.
An $\operatorname{argmax}$ operation converts these to discrete
indices that select the next beacon interval from
$\mathcal{T} = \{0.1, 0.2, 0.5, 1.0\}$~s and the transmission
power offset from
$\mathcal{P} = \{-10, 0, +10, +20\}$~dBm, respectively.  The
entire inference executes in a single forward pass requiring no
iterative optimization at run time.

\textbf{ETSI Compliance Gate.}  The raw TinyMLP outputs are passed
through a hard clamp that enforces
$T_{\text{GenCam}} \in [100\,\text{ms},\, 1000\,\text{ms}]$ before
the timer is re-armed.  Since all four values in $\mathcal{T}$ already
lie within this range, the clamp acts as a safety guard against
implementation faults or adversarial inputs rather than a routine
correction step.  As noted above, any ETSI event trigger issued during
the interval overrides the timer and causes an immediate CAM
regardless of the scheduled $T_{\text{GenCam}}^{\mathrm{AI}}$ value.

\textbf{Learning Pipeline.}  The TinyMLP is trained entirely offline
through Behavior Cloning.  An \emph{oracle} is first constructed by
exhaustive grid-search over the $|\mathcal{T}| \times |\mathcal{P}| = 16$
action combinations for each sampled vehicle state; the oracle
selects the action that minimizes a weighted joint cost combining
AoI and CBR (detailed in Section~\ref{sec.prop}).  The resulting
state--action pairs form the Behavior Cloning dataset.  The TinyMLP
is trained on this dataset with a cross-entropy loss, converted to
ONNX format, post-training quantized to INT8, and exported as a
C header file (\texttt{tinymlp\_aidcc.h}) for on-board deployment.

\textbf{Evaluation Metrics.}  Performance is assessed using five
complementary metrics: (M1)~average Age of Information
($\overline{\text{AoI}}$, in ms), measuring information freshness
across all neighbor pairs; (M2)~Channel Busy Ratio (CBR), measuring
medium utilization; (M3)~Packet Delivery Ratio (PDR), measuring
reception success within 100~m; (M4)~energy efficiency (mJ/km),
measuring distance-normalized transmission energy; and
(M5)~ETSI compliance rate, measuring the fraction of CAM intervals
that respect the standard bounds.  Detailed definitions and
measurement procedures are given in Section~\ref{sec.s}.


% ============================================================
\section{Proposed Scheme}
\label{sec.prop}
% ============================================================

\subsection{State, Action, and Joint Cost Formulation}
\label{subsec:1}

\textbf{Definition 1 (Vehicle State).}  The vehicle state at decision
epoch $k$ is defined as the five-dimensional real-valued vector
\begin{equation}
  \mathbf{s}_k =
  \bigl(v_k,\; a_k,\; N_{\text{est},k},\; \text{CBR}_k,\; \text{AoI}_k\bigr)
  \;\in\; \mathbb{R}^{5},
  \label{eq:state}
\end{equation}
where $v_k$~(m/s) is the current longitudinal speed,
$a_k$~(m/s$^2$) is the longitudinal acceleration,
$N_{\text{est},k}$ is the estimated number of active neighbors within
the communication range (derived from the set of CAMs received within
the preceding 1~s window), $\text{CBR}_k \in [0,1]$ is the
locally measured Channel Busy Ratio over the most recent 100~ms
sensing period, and $\text{AoI}_k$~(ms) is the average Age of
Information computed over all currently observed neighbors.

The AoI of vehicle $j$ as observed by vehicle $i$ at time $t$ is
defined as
\begin{equation}
  \Delta_{ij}(t) = t_{\text{rx}} - t_{\text{gen}},
  \label{eq:aoi_def}
\end{equation}
where $t_{\text{rx}}$ is the wall-clock time at which vehicle $i$ most
recently received a CAM from vehicle $j$, and $t_{\text{gen}}$ is the
generation timestamp embedded in that CAM's payload.  The network-wide
average AoI is
\begin{equation}
  \overline{\text{AoI}}_k =
  \frac{1}{|{\mathcal{N}_i}|}
  \sum_{j \in \mathcal{N}_i} \Delta_{ij}(t_k),
  \label{eq:aoi_avg}
\end{equation}
where $\mathcal{N}_i$ denotes the set of neighbors currently visible to
vehicle $i$.

\textbf{Action Space.}  The joint action selected at epoch $k$ is
\begin{equation}
  (T_k,\, p_k) \;\in\; \mathcal{T} \times \mathcal{P},
  \label{eq:action_space}
\end{equation}
where the discrete beacon interval set is
$\mathcal{T} = \{0.1,\, 0.2,\, 0.5,\, 1.0\}$~s
and the discrete transmission power set is
$\mathcal{P} = \{-10,\, 0,\, +10,\, +20\}$~dBm.
The action space therefore contains $|\mathcal{T}| \times |\mathcal{P}|
= 16$ elements, which is small enough for exhaustive evaluation yet
sufficiently expressive to cover the full ETSI-compliant operating range.

\textbf{Joint Cost Function.}  For oracle construction and ablation
analysis, we employ the following weighted joint cost
\begin{equation}
  J(T, p \mid \mathbf{s}) =
  \alpha \cdot \overline{\text{AoI}}
  \;+\;
  \beta \cdot \text{CBR}
  \;+\;
  \gamma \cdot \mathbb{1}\bigl[\text{ETSI violation}\bigr] \cdot M,
  \label{eq:cost}
\end{equation}
where $\overline{\text{AoI}}$ is evaluated in milliseconds (normalised
to the range $[0,1]$ during training), $\text{CBR} \in [0,1]$, and
the indicator $\mathbb{1}[\text{ETSI violation}]$ equals 1 if the
proposed $T$ falls outside $[T_{\min}, T_{\max}]$.  The penalty
constant $M$ is set to a large value (e.g.,\ $M = 10^4$) to make
ETSI violations strictly infeasible in the oracle solution.
The default weight vector is $(\alpha, \beta, \gamma) = (0.5, 0.5, 1.0)$;
sensitivity analysis over $\alpha$ is reported in Section~\ref{sec.s}.

\subsection{Oracle Generation via Exhaustive Grid Search}
\label{subsec:2}

For each sampled vehicle state $\mathbf{s}$, the oracle evaluates all
16 candidate actions in $\mathcal{T} \times \mathcal{P}$ by running a
200~ms forward simulation.  The action that achieves the minimum cost
$J^{*}(\mathbf{s})$ is assigned as the oracle label:
\begin{equation}
  (T^{*}, p^{*}) = \underset{(T,p)\,\in\,\mathcal{T}\times\mathcal{P}}
  {\operatorname{arg\,min}}\;
  J(T, p \mid \mathbf{s}).
  \label{eq:oracle}
\end{equation}

\textbf{Proposition 1 (Myopic Optimality).}
\textit{For any fixed vehicle state $\mathbf{s}$ and evaluation
horizon of 200~ms, the oracle action $(T^{*}, p^{*})$ defined
in~\eqref{eq:oracle} is globally optimal over $\mathcal{T} \times
\mathcal{P}$, because the grid is exhaustively enumerated.}

\noindent
We stress that this optimality is \emph{myopic}: the oracle minimises
the single-step cost within the 200~ms window and does not account for
long-horizon trajectory effects.  This is an acceptable approximation
because the CAM scheduling decision at epoch $k$ has limited
temporal influence beyond two to three beacon intervals, and the
resulting Behavior Cloning policy is evaluated over its long-run
average rather than its per-step optimality.

The oracle label generation procedure is summarised in
Algorithm~\ref{alg:oracle}.

\begin{algorithm}[!t]
  \caption{Oracle Label Generation via Exhaustive Grid Search}
  \label{alg:oracle}
  \begin{algorithmic}[1]
    \REQUIRE State samples $\mathcal{D} = \{\mathbf{s}^{(i)}\}_{i=1}^{n}$,
             action grid $\mathcal{T} \times \mathcal{P}$,
             horizon $H = 200$~ms,
             cost weights $\alpha, \beta, \gamma$, penalty $M$
    \ENSURE  Labelled dataset
             $\mathcal{L} = \{(\mathbf{s}^{(i)}, T^{*(i)}, p^{*(i)})\}_{i=1}^{n}$
    \FOR{each state sample $\mathbf{s}^{(i)} \in \mathcal{D}$}
      \STATE $J^{*} \leftarrow +\infty$;
             $(T^{*}, p^{*}) \leftarrow \text{null}$
      \FOR{each $(T, p) \in \mathcal{T} \times \mathcal{P}$}
        \STATE Apply $(T, p)$ to the simulator initialised at
               $\mathbf{s}^{(i)}$
        \STATE Run simulation for $H = 200$~ms; record
               $\overline{\text{AoI}}$, $\text{CBR}$,
               and ETSI violation flag
        \STATE Compute $J(T, p \mid \mathbf{s}^{(i)})$
               using~\eqref{eq:cost}
        \IF{$J(T, p \mid \mathbf{s}^{(i)}) < J^{*}$}
          \STATE $J^{*} \leftarrow J(T, p \mid \mathbf{s}^{(i)})$;
                 $(T^{*}, p^{*}) \leftarrow (T, p)$
        \ENDIF
      \ENDFOR
      \STATE Append $(\mathbf{s}^{(i)}, T^{*}, p^{*})$ to $\mathcal{L}$
    \ENDFOR
    \RETURN $\mathcal{L}$
  \end{algorithmic}
\end{algorithm}

\noindent
The dataset $\mathcal{L}$ is generated offline from libsumo
simulations covering the urban-grid scenario at three vehicle
densities, yielding at least 10,000 labelled state--action pairs
before train/validation/test splitting.

\subsection{TinyMLP Architecture and Behavior Cloning}
\label{subsec:3}

\textbf{Architecture.}
The TinyMLP maps the 5D state vector to a pair of four-class softmax
logit blocks.  The layer-by-layer structure is given in
Table~\ref{tab:arch}.

\begin{table}[!t]
  \centering
  \caption{TinyMLP Layer-by-Layer Architecture}
  \label{tab:arch}
  \renewcommand{\arraystretch}{1.2}
  \begin{tabular}{clcc}
    \toprule
    \textbf{Layer} & \textbf{Operation} & \textbf{Output Shape} & \textbf{Params} \\
    \midrule
    0 & Input              & $(5)$    & ---   \\
    1 & FC$(5 \to 32)$ + ReLU  & $(32)$   & $5{\times}32 + 32 = 192$   \\
    2 & FC$(32 \to 32)$ + ReLU & $(32)$   & $32{\times}32 + 32 = 1056$ \\
    3 & FC$(32 \to 8)$         & $(8)$    & $32{\times}8 + 8 = 264$    \\
    4 & Reshape$(4,\,2)$       & $(4, 2)$ & ---   \\
    5 & $\operatorname{argmax}(\cdot, \text{axis}=0)$ & $(2)$ & --- \\
    \bottomrule
  \end{tabular}
\end{table}

\noindent
The total parameter count is
\begin{equation}
  P_{\text{total}}
  = (5 \times 32 + 32)
  + (32 \times 32 + 32)
  + (32 \times 8 + 8)
  = 192 + 1056 + 264
  = 1{,}448.
  \label{eq:param_count}
\end{equation}

Column 2 of the Reshape layer splits the eight output logits into two
groups of four: the first four correspond to the $T_{\text{GenCam}}$
index and the second four to the $p_{\text{tx}}$ index.
At inference time, $\operatorname{argmax}$ is applied independently
to each group, yielding the selected action indices.

\textbf{Behavior Cloning Loss.}
Given a labelled example $(\mathbf{s}, T^{*}, p^{*})$, let
$\hat{\mathbf{l}}_{T} \in \mathbb{R}^4$ and
$\hat{\mathbf{l}}_{p} \in \mathbb{R}^4$ denote the predicted logit
vectors for the beacon interval and power, respectively.  The
training objective is
\begin{equation}
  \mathcal{L}_{\text{BC}} =
  \mathrm{CE}\!\left(\hat{\mathbf{l}}_{T},\, T^{*}\right)
  +
  \mathrm{CE}\!\left(\hat{\mathbf{l}}_{p},\, p^{*}\right),
  \label{eq:bc_loss}
\end{equation}
where $\mathrm{CE}(\cdot, \cdot)$ denotes the standard
cross-entropy loss and $T^{*}$, $p^{*}$ are interpreted as
one-hot class indices.  The two losses are summed rather than
averaged to give equal weight to both output heads; no additional
task-balancing hyperparameter is required.

\textbf{Training Setup.}
The model is trained with the Adam optimiser at an initial learning
rate of $10^{-3}$, a mini-batch size of 64, and for a maximum of
100 epochs.  Training is performed on a standard CPU; given the
model size of 1,448 parameters, a single epoch completes in under
one second.  The labelled dataset is split 70/15/15 into
training, validation, and test sets.  Early stopping is applied
if the validation loss does not improve for 10 consecutive epochs,
preventing overfitting on the relatively compact dataset.  The
final model is selected by the lowest validation cross-entropy.

\subsection{ETSI Compliance and MCU Deployment}
\label{subsec:4}

\textbf{ETSI Compliance Guarantees.}
Two mechanisms ensure that the AI-DCC never violates the ETSI
EN~302~637-2 standard.  First, the beacon interval output is hard-clamped:
\begin{equation}
  T_{\text{GenCam}}^{\text{deploy}}
  = \operatorname{clamp}\!\left(
      T_{\text{GenCam}}^{\mathrm{AI}},\;
      T_{\min},\; T_{\max}
    \right)
  = \operatorname{clamp}\!\left(
      T_{\text{GenCam}}^{\mathrm{AI}},\;
      0.1\,\text{s},\; 1.0\,\text{s}
    \right).
  \label{eq:clamp}
\end{equation}
Because all members of $\mathcal{T}$ are already within
$[0.1\,\text{s},\, 1.0\,\text{s}]$, this clamp is redundant in the
nominal operating path but acts as a fail-safe against
implementation-level faults.  Second, and more importantly,
the ETSI event-driven triggers (Delta Heading $> 4^{\circ}$,
Delta Position $> 4$~m, Delta Speed $> 0.5$~m/s) are evaluated
continuously and always supersede the AI timer: when any trigger
fires, a CAM is generated immediately without consulting the
TinyMLP, and the AI timer is re-armed using the most recent
$T_{\text{GenCam}}^{\text{deploy}}$ value.  This priority ordering
guarantees that the AI layer never suppresses a safety-critical
event-triggered beacon~\cite{Bhattacharyya2024,Iliopoulos2025}.

\textbf{Post-Training INT8 Quantization.}
After training, the floating-point model is converted to ONNX
format and then quantized to INT8 precision using Post-Training
Quantization (PTQ) calibrated on the validation set.  The
quantized model is exported to a flat C array representation
suitable for static linking with CMSIS-NN.  INT8 inference on
ARM Cortex-M hardware is hardware-accelerated through the SIMD
dot-product instructions available in the Cortex-M4 DSP
extension, which processes four INT8 multiply-accumulate
operations per clock cycle.

\textbf{MCU Deployment Profile.}
The reference deployment target is an STM32F407 microcontroller
(ARM Cortex-M4, 168~MHz clock, 1~MB Flash, 192~KB SRAM).
Based on publicly available CMSIS-NN benchmarks, a fully connected
INT8 layer of dimension 32$\times$32 executes in approximately
50--200~cycles on this platform.  Aggregating over the four
weight layers of the TinyMLP, the total inference cost is
estimated at $\sim$500~cycles, corresponding to approximately
$3\,\mu$s at 168~MHz.  This is five orders of magnitude below
the 100~ms minimum CAM interval, leaving ample computational
headroom for the remainder of the CAM service processing.

The memory footprint breaks down as follows.  The INT8 weight
arrays occupy approximately 1.4~KB of Flash.  The largest
intermediate activation buffer (the first hidden layer, 32
neurons) requires 32 bytes; the full stack of activations
across all layers occupies approximately 0.5~KB, giving a
combined on-chip memory requirement of
\begin{equation}
  M_{\text{total}}
  \approx 1.4\,\text{KB (weights)}
  + 0.5\,\text{KB (activations)}
  \approx 2\,\text{KB}.
  \label{eq:memory}
\end{equation}
This is less than 2\% of the STM32F407's 192~KB SRAM, leaving
the remainder entirely available for the ETSI CAM service buffers,
DCC state variables, and the operating system.

\textbf{Deployment Artifact.}
The quantized model is packaged as a single C header file
\texttt{tinymlp\_aidcc.h}, which contains the weight arrays,
bias arrays, layer dimensions, and a self-contained
\texttt{aiDCC\_infer(float *state, uint8\_t *T\_idx,
uint8\_t *p\_idx)} function callable from any ETSI CAM service
implementation without external library dependencies beyond
CMSIS-NN~\cite{Zila2026}.


% ============================================================
\section{Performance Evaluation}
\label{sec.s}
% ============================================================

\subsection{Simulation Environment and Datasets}
\label{sec.s0}

All experiments are conducted using \textbf{libsumo} (the Python
binding of the SUMO traffic simulator) for vehicle mobility generation
and \textbf{SumoNetSim~1.1.5} for the IEEE~802.11p MAC/PHY channel
simulation stack.  The map scenario is \texttt{urban\_grid}, which
models a rectangular grid of road segments totalling 14.4~km of
drivable length, organized as two-lane bidirectional roads over three
blocks.  The vehicle population parameter is set to
\texttt{n\_vehicles}~=~50 (ETSI compatibility nominal); under the
dynamic trip-generation model the actual average number of concurrently
active vehicles is approximately 88, yielding an equivalent vehicle
density of approximately \textbf{3.06~veh/(km$\cdot$lane)}.  Each
simulation runs for \texttt{duration\_steps}~=~3600 seconds of
simulated time.  To assess result reproducibility, three independent
random seeds---\{42, 123, 456\}---are used for each method, a set that
provides sufficient statistical determinism for ETSI compliance
verification.  Seed-wise coefficient of variation remained below 1.1\%
for both AoI and CBR across all four baselines, confirming
reproducibility of the reported means.

The physical layer follows the ITS-G5 specification at 5.9~GHz with a
10~MHz channel bandwidth and IEEE~802.11p OFDM modulation.  Small-scale
fading is modelled by a Nakagami-$m$ distribution and large-scale
propagation by a distance-based path-loss with exponent
$\alpha = 2.0$, representative of near-line-of-sight urban road
segments (consistent with the model presented in
Section~\ref{sec.net}).  Raw simulation output is stored in per-method
CSV files under \texttt{paper/data/main\_BL-*\_urban.csv}; aggregated
cross-method statistics are stored in
\texttt{paper/data/main\_combined\_urban.csv}.

\subsection{Baselines and Metrics}
\label{sec.s_setup}

\textbf{Baselines.}
Four baseline controllers are evaluated, spanning the design space
from the standard ETSI Reactive algorithm to a high-rate fixed-interval
reference:

\begin{itemize}[leftmargin=2em]
  \item[\textbf{BL-A}] \textbf{ETSI DCC Reactive}
    (ETSI EN~302~571 \S8.1)~\cite{Bhattacharyya2024}: the
    standard three-state CBR-threshold DCC with fixed transmission
    power $p_{\text{tx}} = +20$~dBm.  This is the primary regulatory
    baseline.

  \item[\textbf{BL-B}] \textbf{Simplified Adaptive DCC}: a
    proportional CBR-tracking controller with
    CBR$_{\text{target}} = 0.6$, update step $\Delta T = 0.05$~s,
    and smoothing factor $\lambda_s = 0.5$.  Power is fixed at
    $+20$~dBm.  This baseline represents the practical
    ``Adaptive'' DCC variant referenced in idea\_spec \S5.3 and
    \S10.1.

  \item[\textbf{BL-C}] \textbf{Bhattacharyya2024 Variable Beacon}~\cite{Bhattacharyya2024}:
    a CBR-plus-neighbor-count heuristic that adjusts
    $T_{\text{GenCam}}$ based on both channel occupancy and locally
    estimated neighbor density, without ML inference.

  \item[\textbf{BL-D}] \textbf{Fixed 10~Hz}: a constant-rate
    reference operating at the shortest permissible interval
    ($T_{\text{GenCam}} = 100$~ms, $p_{\text{tx}} = +20$~dBm),
    which maximises information freshness at the cost of channel
    load.
\end{itemize}

\textbf{Metrics.}
Performance is quantified by five complementary metrics (M1--M5),
whose formal definitions are given in Section~\ref{sec.net}; here
we briefly restate them for completeness:
\begin{itemize}[leftmargin=2em]
  \item[M1] \textbf{Mean AoI} ($\overline{\text{AoI}}$, ms) ---
    average Age of Information across all observable neighbor pairs,
    measuring information freshness.
  \item[M2] \textbf{CBR} --- Channel Busy Ratio, measuring medium
    utilization.
  \item[M3] \textbf{PDR} (\%) --- Packet Delivery Ratio within 100~m,
    measuring reception success.
  \item[M4] \textbf{Energy Efficiency} (dimensionless proxy) ---
    distance-normalized transmission-energy surrogate.
  \item[M5] \textbf{ETSI Compliance} (\%) --- fraction of CAM
    generation intervals that respect $[T_{\min}, T_{\max}]$.
\end{itemize}

\noindent
Fair comparison across methods is ensured by using the identical set
of three seeds, the same \texttt{.sumocfg} network asset, and the
same warm-up procedure that excludes the initial simulation transient
from all reported statistics.

\subsection{Baseline Comparison Results}
\label{sec.s_baseline}

Table~\ref{tab:baseline_main} summarises the mean performance of
the four baselines over the three random seeds in the
\texttt{urban\_grid} scenario.

\begin{table*}[!t]
  \centering
  \caption{Phase-2 main comparison across four baselines under the
    \texttt{urban\_grid} scenario (mean over three random seeds;
    per-method seed CV\,$<$\,1.1\%). Proposed evaluation is pending
    the completion of the Behavior-Cloning training pipeline.}
  \label{tab:baseline_main}
  \renewcommand{\arraystretch}{1.25}
  \begin{tabular}{lcccccc}
    \toprule
    \textbf{Method}
      & \textbf{AoI (ms)}
      & \textbf{CBR}
      & \textbf{PDR (\%)}
      & \textbf{Energy Eff.}
      & \textbf{ETSI Comp. (\%)}
      & \textbf{\# CAM Events} \\
    \midrule
    BL-A (ETSI Reactive)
      & \underline{393.81}
      & \underline{0.5065}
      & 100.00
      & \underline{5.707}
      & 100.00
      & 1{,}676{,}693 \\
    BL-B (Simplified Adaptive)
      & 380.96
      & 0.5115
      & 100.00
      & 5.708
      & 100.00
      & 1{,}677{,}008 \\
    BL-C (Variable Beacon~\cite{Bhattacharyya2024})
      & 501.41
      & \textbf{0.4098}
      & 100.00
      & 4.555
      & 100.00
      & 1{,}338{,}457 \\
    BL-D (Fixed 10~Hz)
      & \textbf{321.96}
      & 0.6208
      & 100.00
      & \textbf{6.888}
      & 100.00
      & 2{,}023{,}791 \\
    \midrule
    Proposed (TinyMLP-AI-DCC)
      & \multicolumn{1}{c}{---}
      & \multicolumn{1}{c}{---}
      & \multicolumn{1}{c}{---}
      & \multicolumn{1}{c}{---}
      & \multicolumn{1}{c}{---}
      & \multicolumn{1}{c}{Pending} \\
    \bottomrule
  \end{tabular}
  \smallskip\\
  \small \textbf{Bold} = best in column; \underline{Underline} = second-best in column.
  AoI and CBR are the primary trade-off metrics; all methods achieve perfect PDR and ETSI compliance.
\end{table*}

The four baselines collectively reveal the fundamental AoI--CBR
trade-off inherent in the \texttt{urban\_grid} scenario.
\textbf{BL-D} (Fixed~10~Hz) achieves the lowest mean AoI of
\textbf{321.96}~ms by transmitting at the maximum permissible rate,
providing the freshest cooperative awareness among all evaluated
methods.  However, this aggressive beaconing strategy produces the
highest CBR of 0.6208, which is already at the boundary of the ETSI
Restricted zone (CBR $> 0.60$), indicating that Fixed~10~Hz
operation places the channel under near-overload conditions.  In
high-density or mixed-traffic scenarios this CBR level would be
unsustainable, as any additional vehicular load would push the network
into congestion collapse.

\textbf{BL-C} (Variable Beacon) occupies the opposite extreme of the
trade-off curve.  By aggressively reducing the beacon interval in
response to CBR and neighbor-count signals, BL-C achieves the lowest
channel utilization of \textbf{0.4098}---well within the ETSI Relaxed
zone---but at the cost of the worst mean AoI of 501.41~ms among all
baselines.  The 55.8\% increase in AoI relative to BL-D suggests that
the variable-beacon heuristic over-reduces the CAM rate, leaving
neighbor state information unacceptably stale for cooperative
perception applications.  The correspondingly lower energy efficiency
(4.555) further indicates that BL-C's under-utilization of the channel
does not translate into improved system performance.

The two ETSI-aligned variants---\textbf{BL-A} (Reactive) and
\textbf{BL-B} (Simplified Adaptive)---occupy the middle of the
trade-off space with comparable, moderate AoI and CBR values.
Notably, BL-B improves mean AoI by 12.9~ms ($-3.3\%$) relative to
BL-A while incurring only a marginal 0.5~percentage-point increase in
CBR (0.5115 vs.\ 0.5065).  This result confirms that the
proportional CBR-tracking loop of the Simplified Adaptive controller
provides slightly more timely information updates on average, under
essentially equivalent channel load---a modest but consistent
advantage of adaptive over purely reactive DCC.

The distribution of operating points described above directly
motivates the proposed TinyMLP-AI-DCC design.  The four baselines
define a frontier in the AoI--CBR plane that no single heuristic can
simultaneously dominate: BL-D minimises AoI but saturates the channel,
while BL-C minimises CBR but sacrifices freshness.  The proposed
controller targets a \emph{Pareto-dominant} operating point---driving
AoI toward the BL-D level while maintaining CBR in the BL-A/B range
(approximately 0.50--0.51)---a region that lies strictly inside the
convex hull of the baseline operating points and is unreachable by any
of the four static-policy baselines.  This target aligns precisely
with the minimization of the joint cost $J(T, p)$ defined in
Section~\ref{sec.prop}, Eq.~\eqref{eq:cost}, where the
$\alpha$-weighted AoI term and the $\beta$-weighted CBR term enforce
exactly this kind of balanced trade-off.

\subsection{Pending Evaluation Items}
\label{sec.s_pending}

Owing to the staged development of the TinyMLP-AI-DCC pipeline, the
proposed-vs-baseline comparison, ablation analysis, and sensitivity
sweeps are deferred to the subsequent submission revision; the current
baseline study isolates the operating-point gap that motivates the
proposed joint controller.  The following evaluation items remain
pending:

\begin{itemize}[leftmargin=2em]
  \item \textbf{Proposed (TinyMLP-AI-DCC) end-to-end evaluation} ---
    full five-metric comparison against the four baselines upon
    completion of the Behavior-Cloning training pipeline.
  \item \textbf{Ablation study} --- ABL-1 (rate-only control, power
    fixed) and ABL-2 (no AoI input, 4D state vector) to isolate
    the contribution of joint control and AoI awareness respectively.
  \item \textbf{Sensitivity analysis} --- SA1 ($\alpha/\beta$ cost
    weight ratio), SA2 (grid size / road topology), SA3 (CBR
    target threshold), SA4 (TinyMLP hidden layer width).
  \item \textbf{Additional scenarios} --- Scenario S2 (highway
    mobility pattern) and a vehicle density sweep over
    10, 20, and 30~veh/(km$\cdot$lane) to assess robustness across
    traffic conditions.
\end{itemize}

% ============================================================
\section{Conclusion}
\label{sec.c}
% ============================================================

This paper addressed the fundamental tension between information
freshness and channel congestion in ETSI CAM-based vehicular IoT
networks.  Standard DCC algorithms, being reactive and CBR-centric,
cannot jointly optimize the Age of Information and channel utilization,
nor do they leverage the kinematic context of individual vehicles.
We proposed \textbf{TinyMLP-AI-DCC}, a two-layer MLP with 1,448
parameters that accepts a five-dimensional vehicle state vector
(speed, acceleration, estimated neighbor count, CBR, and AoI) and
outputs a joint action $(T_{\text{GenCam}}, p_{\text{tx}})$ in a
single forward pass.  The model is trained offline via Behavior
Cloning against a myopically optimal oracle generated by exhaustive
16-point grid search, avoiding the convergence instability of deep
reinforcement learning.  Post-training INT8 quantization reduces the
deployment footprint to under 4~KB of MCU memory---fitting
comfortably within the 192~KB SRAM of an STM32F407 (ARM Cortex-M4)
and completing inference in under 1~ms at 80~MHz.

The baseline comparison study reported in this work provides
quantitative evidence for the AoI--CBR Pareto gap that motivates the
proposed design.  Across the \texttt{urban\_grid} scenario (three
random seeds, $\sim$88 concurrent vehicles, density $\approx$
3.06~veh/(km$\cdot$lane)), the four baseline methods reveal a clear
frontier: Fixed~10~Hz (BL-D) achieves the lowest AoI of 321.96~ms
but pushes CBR to 0.6208 (near-overload); Variable Beacon (BL-C)
achieves the lowest CBR of 0.4098 but at an AoI cost of 501.41~ms;
while the two ETSI-standard variants (BL-A/B) occupy a moderate but
unoptimized middle ground.  No existing baseline simultaneously
achieves low AoI \emph{and} moderate CBR, confirming the existence of
an open Pareto-dominant operating region that the TinyMLP-AI-DCC
controller is designed to target, in direct correspondence with the
joint cost $J(T, p)$ defined in Section~\ref{sec.prop}.

\textbf{Limitations.}  The present version of this paper has two
principal limitations that must be acknowledged.  First, the
end-to-end evaluation of the proposed TinyMLP-AI-DCC controller---its
training, BC policy convergence, and five-metric comparison against
the four baselines---has not yet been completed, and all entries in
Table~\ref{tab:baseline_main} for the Proposed method remain pending.
Consequently, the performance claims in the abstract (``up to 35\%
AoI reduction'') are projections based on the oracle performance
analysis rather than empirical end-to-end results.  Second, the MCU
resource figures (inference latency and memory footprint) are
simulation estimates derived from publicly available CMSIS-NN
benchmarks; no physical measurement on actual STM32F407 hardware has
been performed in this version.

\textbf{Future Work.}
The immediate next steps are: (i)~completion of the TinyMLP Behavior
Cloning training pipeline and end-to-end simulation evaluation;
(ii)~ablation studies (ABL-1: rate-only, ABL-2: no-AoI input) and
sensitivity sweeps (SA1--SA4) as detailed in
Section~\ref{sec.s_pending}; (iii)~extension to the highway mobility
scenario~(S2) and a density sweep across 10--30~veh/(km$\cdot$lane)
to characterise the controller's generalization envelope; and
(iv)~physical deployment on an actual STM32F407 board to obtain
measured latency and memory consumption data, replacing the current
simulation-based estimates and providing a concrete, reproducible MCU
benchmark.

% ============================================================
\begin{thebibliography}{00}
% ********** Related Work: AI-driven MAC Protocols (Direct competition) **********
% (10 papers)
\bibitem{Wu2025} Chien-Min Wu, Cheng-Tai Tsai, Cheng-Chun Hou, \emph{et al.}, ``Emergency Message Broadcast Mechanism in Vehicular Ad-Hoc Networks Based on Reinforcement Learning With Contention Estimation,'' \emph{IEEE Transactions on Intelligent Vehicles}, 2025. doi: 10.1109/TIV.2024.3418778.
\bibitem{Iqbal2025} Adeel Iqbal, A. Nauman, Tahir Khurshaid, ``Lightweight Reinforcement Learning for Priority-Aware Spectrum Management in Vehicular IoT Networks,'' \emph{MDPI Sensors}, 2025. doi: 10.3390/s25216777.
\bibitem{Louvros2026} Spyridon Louvros, Anupkumar Pandey, B. Shah, Yashesh Buch, ``Federated Edge Learning MAC Cross-Layer Optimization in 6G Networks : The Case for Radial Basis Functions Optimization,'' \emph{2026 Panhellenic Conference on Electronics & Telecommunications (PACET)}, 2026. doi: 10.1109/PACET68758.2026.11498274.
\bibitem{Li2025a} Zihan Li, Ping Wang, Yamin Shen, Song Li, ``Reinforcement Learning-Based Resource Allocation Scheme of NR-V2X Sidelink for Joint Communication and Sensing,'' \emph{MDPI Sensors}, 2025. doi: 10.3390/s25020302.
\bibitem{Wang2024} Ping Wang, C. Zheng, Chenggen Pu, Zhiyi Tian, ``Dynamic Multi-Channel Access in Heterogeneous Industrial Wireless Sensor Networks: A MADRL Approach,'' \emph{ACM Cloud and Autonomic Computing Conference}, 2024. doi: 10.1109/CAC63892.2024.10865069.
\bibitem{Ibrahim2025} Amin S. Ibrahim, Ahmed M. Abbas, Saeed Mohsen, ``A novel adaptive IEEE 802.11ah MAC protocol using machine learning for data collision reduction in smart cities IoT networks,'' \emph{Wireless networks}, 2025. doi: 10.1007/s11276-025-04059-2.
\bibitem{Ibrahim2026} Amin S. Ibrahim, Ahmed M. Abbas, Saeedi Mohsen, ``Correction: A novel adaptive IEEE 802.11ah MAC protocol using machine learning for data collision reduction in smart cities IoT networks,'' \emph{Wireless networks}, 2026. doi: 10.1007/s11276-025-04078-z.
\bibitem{Zila2026} Amine Zila, Y. Mouzouna, A. Ouchatti, Ikram Daanoune, ``Edge AI and TinyML for Enhancing MAC Protocols: A New Paradigm for Wireless Sensor Networks in IIoT,'' \emph{International Journal of Communication Systems}, 2026. doi: 10.1002/dac.70403.
\bibitem{Ni2024} Yang Ni, Danny Abraham, Mariam Issa, \emph{et al.}, ``Dynamic MAC Protocol for Wireless Spectrum Sharing via Hyperdimensional Self-Learning,'' \emph{IEEE Access}, 2024. doi: 10.1109/ACCESS.2024.3464868.
\bibitem{PV2025} I. P., T. Venkatesh, ``RIS-Assisted Full Duplex MAC Protocol for NR-V2X with ML-Driven QoS Prediction,'' \emph{2025 IEEE Future Networks World Forum (FNWF)}, 2025. doi: 10.1109/FNWF66845.2025.11317280.

% ********** Related Work: Lightweight ML for Vehicular Networks **********
% (13 papers)
\bibitem{Li2025} Junjun Li, Yanyan Ma, Jiahui Bai, \emph{et al.}, ``A Lightweight Intrusion Detection System with Dynamic Feature Fusion Federated Learning for Vehicular Network Security,'' \emph{MDPI Sensors}, 2025. doi: 10.3390/s25154622.
\bibitem{Yan2025} Na Yan, Senior Member Ieee Kezhi Wang, Kangda Zhi, \emph{et al.}, ``Secure and Private Over-the-Air Federated Learning: Biased and Unbiased Aggregation Design,'' \emph{IEEE Transactions on Wireless Communications}, 2025. doi: 10.1109/TWC.2025.3550159.
\bibitem{Liu2024} Yiming Liu, Qiang Wang, Jiaao Chen, Wenqi Zhang, Chen Sun, ``Federated Multi-Agent Deep Reinforcement Learning Approach for Resource Allocation in Platoon-Based NR-V2X,'' \emph{IEEE Vehicular Technology Conference}, 2024. doi: 10.1109/VTC2024-Spring62846.2024.10683498.
\bibitem{Wu2024} Maoqiang Wu, Guoliang Cheng, Dongdong Ye, \emph{et al.}, ``Federated Split Learning With Data and Label Privacy Preservation in Vehicular Networks,'' \emph{IEEE Transactions on Vehicular Technology}, 2024. doi: 10.1109/TVT.2023.3304176.
\bibitem{Zhang2025} Jianquan Zhang, Fangting Huang, Shuqing Zhu, Xiao Xiao, ``A Resource Allocation Strategy in Internet of Vehicles Based on Multi-Task Federated Learning and Incentive Mechanism,'' \emph{IEEE transactions on intelligent transportation systems (Print)}, 2025. doi: 10.1109/TITS.2025.3528969.
\bibitem{Yang2024a} Lu Yang, Songtao Guo, Chen-Khong Tham, \emph{et al.}, ``PreM-FedIoV: A Novel Federated Reinforcement Learning Framework for Predictive Maintenance in IoV,'' \emph{IEEE Transactions on Mobile Computing}, 2024. doi: 10.1109/TMC.2024.3404249.
\bibitem{Xu2024} Mengfan Xu, Yaguang Lin, ``FedSteg: Coverless Steganography‐Based Privacy‐Preserving Decentralized Federated Learning,'' \emph{IEEJ TRANSACTIONS ON ELECTRICAL AND ELECTRONIC ENGINEERING}, 2024. doi: 10.1002/tee.24085.
\bibitem{Elzemity2024} Adel Elzemity, Budi Arief, ``Privacy Threats and Countermeasures in Federated Learning for Internet of Things: A Systematic Review,'' \emph{2024 IEEE International Conferences on Internet of Things (iThings) and IEEE Green Computing & Communications (GreenCom) and IEEE Cyber, Physical & Social Computing (CPSCom) and IEEE Smart Data (SmartData) and IEEE Congress on Cybermatics}, 2024. doi: 10.1109/iThings-GreenCom-CPSCom-SmartData-Cybermatics62450.2024.00072.
\bibitem{Liu2025} Xianhui Liu, Jianle Liu, Yingyao Zhang, Ning Jia, Chenlin Zhu, ``Heterogeneous Federated Learning for Vehicle-to-Everything: Feature Prototype Aggregation and Generative Feedback Mechanism,'' \emph{IEEE Open Journal of Vehicular Technology}, 2025. doi: 10.1109/OJVT.2025.3594030.
\bibitem{Zhang2025b} Zhiwei Zhang, ``A Vehicular Network Intrusion Detection Method Based on Knowledge Distillation and Hybrid Learning,'' \emph{2025 International Conference on Algorithms, Data Mining, and Information Technology (ADMIT)}, 2025. doi: 10.1109/ADMIT67050.2025.11337192.
\bibitem{Liao2025} Gengjian Liao, Yulin Yang, Zhenni Feng, ``Personalized federated learning through self-knowledge distillation in vehicular edge computing,'' \emph{Comput. Networks}, 2025. doi: 10.1016/j.comnet.2025.111406.
\bibitem{Li2025b} Chunlin Li, Yong Zhang, Long Yu, Mengjie Yang, ``Efficient Vehicle Selection and Resource Allocation for Knowledge Distillation-Based Federated Learning in UAV-Assisted VEC,'' \emph{IEEE transactions on intelligent transportation systems (Print)}, 2025. doi: 10.1109/TITS.2025.3525735.
\bibitem{Zhao2026} Pengcheng Zhao, Xingjia Wei, Yong Wang, Xin Li, Shibao Sun, ``AFed-CHKD: An Adversarial Federated Learning-Based Cluster Head Selection and Knowledge Distillation in Vehicular Sensor Network,'' \emph{IEEE Sensors Journal}, 2026. doi: 10.1109/JSEN.2026.3665994.

% ********** Related Work: V2X / NR-V2X Protocols **********
% (10 papers)
\bibitem{Nikooroo2024} Saleh Nikooroo, Juan Estrada-Jimenez, Aurel Machalek, \emph{et al.}, ``Poster: Evaluating the Potential of Mode 2(b) Resource Allocation in NR V2X Sidelink,'' \emph{IEEE Vehicular Networking Conference}, 2024. doi: 10.1109/VNC61989.2024.10575999.
\bibitem{Li2024} Pengfei Li, Xin-Lin Huang, Fei Hu, ``When DSRC Meets With C-V2X: How to Jointly Predict and Select Idle Channels?,'' \emph{IEEE Transactions on Vehicular Technology}, 2024. doi: 10.1109/TVT.2024.3399211.
\bibitem{Gao2024} Ang Gao, Ziqing Zhu, Jiankang Zhang, Wei Liang, Yansu Hu, ``Matching Combined Heterogeneous Multi-Agent Reinforcement Learning for Resource Allocation in NOMA-V2X Networks,'' \emph{IEEE Transactions on Vehicular Technology}, 2024. doi: 10.1109/TVT.2024.3409048.
\bibitem{Narayanasamy2024} Iswarya Narayanasamy, Venkateswari Rajamanickam, ``A Cascaded Multi-Agent Reinforcement Learning-Based Resource Allocation for Cellular-V2X Vehicular Platooning Networks,'' \emph{MDPI Sensors}, 2024. doi: 10.3390/s24175658.
\bibitem{Chen2025} Chen Chen, Wenhua Wang, Ziye Liu, \emph{et al.}, ``RLFN-VRA: Reinforcement Learning-Based Flexible Numerology V2V Resource Allocation for 5G NR V2X Networks,'' \emph{IEEE Transactions on Intelligent Vehicles}, 2025. doi: 10.1109/TIV.2024.3427399.
\bibitem{Mao2024} Yaqi Mao, Xin Yang, Ling Wang, \emph{et al.}, ``A High-Capacity MAC Protocol for UAV-Enhanced RIS-Assisted V2X Architecture in 3-D IoT Traffic,'' \emph{IEEE Internet of Things Journal}, 2024. doi: 10.1109/JIOT.2024.3387997.
\bibitem{Iliopoulos2025} Christos Iliopoulos, Athanasios C. Iossifides, C. Foh, Periklis Chatzimisios, ``IEEE 802.11BD for Next-Generation V2X Communications: From Protocol to Services,'' \emph{IEEE Communications Standards Magazine}, 2025. doi: 10.1109/MCOMSTD.2025.3569015.
\bibitem{Zhang2025a} Zhishuo Zhang, Wen Huang, Ying Huang, \emph{et al.}, ``An Attribute-Based Pre-Authenticated Secure Communication Protocol Enabling Key Protection and Credential Online-Upgrading for 5G NR V2X,'' \emph{IEEE transactions on intelligent transportation systems (Print)}, 2025. doi: 10.1109/TITS.2025.3558978.
\bibitem{Alagumani2024} Sangeetha Alagumani, Umamaheswari Natarajan, ``Q-learning and fuzzy logic multi-tier multi-access edge clustering for 5g v2x communication,'' \emph{Network}, 2024. doi: 10.1080/0954898X.2024.2309947.
\bibitem{Doğan2026} Yasemin Doğan, Ramazan Ünlü, ``Machine Learning for V2X-Enabled Microgrids: A Bibliometric and Thematic Review of Intelligent Energy Management Applications,'' \emph{The Arabian journal for science and engineering}, 2026. doi: 10.1007/s13369-026-11201-5.

% ********** Background: Vehicular Networking, Edge Intelligence **********
% (14 papers)
\bibitem{Rawlley2026} Oshin Rawlley, Shashank Gupta, Jatin Kumar Panwar, Palak Sharma, Shailendra Rathore, ``Asynchronous Deep Reinforcement Learning for Semantic Communication and Digital-Twin Deployment in Transportation Networks,'' \emph{IEEE transactions on intelligent transportation systems (Print)}, 2026. doi: 10.1109/TITS.2025.3593038.
\bibitem{Mianji2025} Elham Mohammadzadeh Mianji, Gabriel-Miro Muntean, Irina Tal, ``Enhancing Vehicular Network Security, Privacy, and Trust Through Reinforcement Learning: A Comprehensive Survey,'' \emph{IEEE transactions on intelligent transportation systems (Print)}, 2025. doi: 10.1109/TITS.2025.3612202.
\bibitem{Yang2024} Liu Yang, Yifei Wei, Zhiyong Feng, Qixun Zhang, Zhu Han, ``Deep Reinforcement Learning-Based Resource Allocation for Integrated Sensing, Communication, and Computation in Vehicular Network,'' \emph{IEEE Transactions on Wireless Communications}, 2024. doi: 10.1109/TWC.2024.3470873.
\bibitem{Geraci2025} Giovanni Geraci, Francesca Meneghello, F. Wilhelmi, \emph{et al.}, ``Wi-Fi: 25 Years and Counting,'' \emph{Proceedings of the IEEE}, 2025. doi: 10.1109/JPROC.2026.3662513.
\bibitem{Xie2025} Xin Xie, Tao Zhong, Heng Wang, ``Scheduling for Maximizing the Information Freshness in Vehicular Edge Computing- Assisted IoT Systems,'' \emph{IEEE transactions on intelligent transportation systems (Print)}, 2025. doi: 10.1109/TITS.2024.3514099.
\bibitem{Chen2026} Xiangyi Chen, Huanlai Xing, Qinnan Zhang, \emph{et al.}, ``Age-of-Information-Aware Hierarchical Collaborative Inference in Digital Twin-Enabled Vehicular Edge Computing,'' \emph{IEEE Transactions on Network Science and Engineering}, 2026. doi: 10.1109/TNSE.2025.3615533.
\bibitem{Liu2024a} Jianhang Liu, Kunlei Xue, Q. Miao, \emph{et al.}, ``MCVCO: Multi-MEC Cooperative Vehicular Computation Offloading,'' \emph{IEEE Transactions on Intelligent Vehicles}, 2024. doi: 10.1109/TIV.2023.3299381.
\bibitem{Huang2025} Mengting Huang, Wenqian Zhang, Guanglin Zhang, ``Joint Service Placement and Task Offloading in Vehicle-Edge-Cloud Collaborative Networks,'' \emph{IEEE transactions on intelligent transportation systems (Print)}, 2025. doi: 10.1109/TITS.2025.3590099.
\bibitem{Zhou2025} Huan Zhou, Wenchao Xia, Gan Zheng, Hongbo Zhu, ``Graph Neural Network-Based Continual Learning for Resource Allocation in Dynamic Wireless Environments,'' \emph{IEEE Transactions on Vehicular Technology}, 2025. doi: 10.1109/TVT.2025.3585146.
\bibitem{Li2024a} Mengfei Li, Wei Wang, Shijie Feng, \emph{et al.}, ``GCN-Enhanced Multi-Agent MAC Protocol for Vehicular Communications,'' \emph{International Conference on Industrial Informatics}, 2024. doi: 10.1109/INDIN58382.2024.10774514.
\bibitem{Tang2024} Zhiqing Tang, Fangyi Mou, Jiong Lou, \emph{et al.}, ``Multi-User Layer-Aware Online Container Migration in Edge-Assisted Vehicular Networks,'' \emph{IEEE/ACM Transactions on Networking}, 2024. doi: 10.1109/TNET.2023.3330255.
\bibitem{Bhattacharyya2024} Sagnik Bhattacharyya, Pankaj Kumar, Sam Darshi, \emph{et al.}, ``Hybrid Relaying Based Cross Layer MAC Protocol Using Variable Beacon for Cooperative Vehicles,'' \emph{IEEE Transactions on Vehicular Technology}, 2024. doi: 10.1109/TVT.2023.3307672.
\bibitem{Atimati2024} Ehinomen Atimati, David H. Crawford, R. Stewart, ``Intelligent Joint Resource Management of Shared Spectrum HetNets,'' \emph{2024 IEEE Future Networks World Forum (FNWF)}, 2024. doi: 10.1109/FNWF63303.2024.11028798.
\bibitem{Alsadie2024} Deafallah Alsadie, ``Artificial Intelligence Techniques for Securing Fog Computing Environments: Trends, Challenges, and Future Directions,'' \emph{IEEE Access}, 2024. doi: 10.1109/ACCESS.2024.3463791.

\end{thebibliography}
% ============================================================

\end{document}
